% ==== Sprache & Mikrotypografie ====
\usepackage{polyglossia} % für mehrsprachige Dokumente
\setmainlanguage{german}
\usepackage{microtype} % für hübschere Typografie (z.B. Ligaturen)
\usepackage{amsmath} % für erweiterte Mathematik-Umgebungen wie z.B. align

% ==== Fonts & Emoji-Fallback (LuaLaTeX) ====
\usepackage{fontspec}
\usepackage{unicode-math}
\usepackage{luacode}

% TikZ für Diagramme und Grafiken
\usepackage{tikz}
\usepackage{tikz-3dplot}
\usetikzlibrary{arrows.meta,calc,arrows,3d,bending,positioning}

% --- Schöner, kompakter Pfeilkopf ---
\tikzset{
  arrowstyle/.style = {-{Stealth[length=4pt, width=4pt]}},
}

% ==== Unicode-Fallbacks ====
% Emoji + allgemeiner Unicode-Fallback
% "Noto Emoji" muss installiert sein
% "DevaVu Sans" muss installiert sein
%
\directlua{
  luaotfload.add_fallback("unifallback", {
    "Noto Emoji:mode=harf;script=DFLT;",
    "DejaVu Sans:mode=harf;script=latn;"
  })
}

% ==== Schriften ====
% Hauptschrift (DIN, mit Fallback)
\setmainfont{HDA DIN Office}[
  BoldFont   = {HDA DIN Office Bold},
  %ItalicFont = {Libertinus Sans Italic},
  RawFeature = {fallback=unifallback}
]

% Sans-Font
\setsansfont{HDA DIN Office}[
  BoldFont   = {HDA DIN Office Bold},
  %ItalicFont = {Libertinus Sans Italic},
  RawFeature = {fallback=unifallback}
]

% Monospace für Code
\setmonofont{DejaVu Sans Mono}[
  Scale      = MatchLowercase,
  RawFeature = {fallback=unifallback}
]

% Mathematische Schrift
\setmathfont{Libertinus Math}

% (Optional hübschere Absätze/Seitenränder – bei Bedarf)
\usepackage[a4paper,margin=2.5cm]{geometry}
\setlength{\parskip}{0.5\baselineskip}
\setlength{\parindent}{0pt}

% Header/Footer
\usepackage{fancyhdr}
\pagestyle{fancy}
\fancyhf{} % alles löschen

% Footer links/rechts
\fancyfoot[C]{\thepage} % Seitenzahl

